\documentclass[11pt]{article}
\usepackage[margin=0.88in]{geometry}
\usepackage{amsmath,amssymb,amsthm,booktabs,graphicx,microtype,enumitem}
\usepackage[hidelinks]{hyperref}

\newtheorem{theorem}{Theorem}
\newtheorem{proposition}[theorem]{Proposition}
\newtheorem{lemma}[theorem]{Lemma}
\newtheorem{corollary}[theorem]{Corollary}
\theoremstyle{definition}
\newtheorem{definition}[theorem]{Definition}
\theoremstyle{remark}
\newtheorem{remark}[theorem]{Remark}
\newcommand{\cl}{\overline}

\title{Left-ideal cores upgrade asymptotic nonexpansiveness}
\author{Candidate substantial partial result for arXiv:2006.15393}
\date{13 August 2026}

\begin{document}
\maketitle

\begin{center}
\fbox{\parbox{0.92\textwidth}{\textbf{Status.}
Candidate substantial partial result, likely valid.  The unrestricted
Question~5.5 remains open here.  We prove two transfer principles under which
every one of the seven fixed-point statements named in that question does
extend to asymptotically nonexpansive actions.  One principle applies to
noncompact amenable semigroups not covered by the compact-semigroup case in
the source.}}
\end{center}

\section{Source and exact question}

The source is Bui Ngoc Muoi and Ngai-Ching Wong, \emph{Fixed point theorems
of various nonexpansive actions of semitopological semigroups on
weakly/weak* compact convex sets}, arXiv:2006.15393v4 \cite{MW}.  It proves
fixed-point results for \emph{super asymptotically nonexpansive} or
\emph{pointwise eventually nonexpansive} actions and closes with the question
reproduced in Figure~\ref{fig:question}.

\begin{figure}[ht]
\centering
\includegraphics[width=0.96\textwidth]{figures/question_crop.png}
\caption{Question~5.5 and the compact-semigroup case, from page 19 of the
official arXiv PDF.}
\label{fig:question}
\end{figure}

\noindent\textbf{Transcription.}
``Do we have similar results as Theorems 3.1, 3.6, 4.1, 4.2, 4.4, and
Propositions 3.11, 3.12 for asymptotically nonexpansive actions?''
The source observes that the answer is affirmative when $S$ is compact and
right reversible.

\section{Definitions}

Let $S$ be a semigroup acting on a subset $K$ of a Banach space.  A
\emph{left ideal} is a nonempty $I\subseteq S$ with $SI\subseteq I$; it is
\emph{supported by} $t\in S$ when $I\subseteq St$.  The action is
\emph{asymptotically nonexpansive} if, for every $x,y\in K$, some left ideal
$I_{x,y}$ satisfies
\begin{equation}
 \|s x-s y\|\leq \|x-y\| \qquad(s\in I_{x,y}).                 \tag{2.1}
\end{equation}
It is \emph{super asymptotically nonexpansive} if, for every $x\in K$ and
$t\in S$, some left ideal $J_{x,t}\subseteq St$ satisfies (2.1) for
\emph{every} $y\in K$.  Omitting the support requirement gives pointwise
eventual nonexpansiveness.  For a locally convex space $(E,Q)$, the
$Q$-versions replace the norm by a specified $q\in Q$ and quantify over $q$.

A \emph{least left ideal} is a left ideal $L$ contained in every left ideal
of $S$.  For an infinite cardinal $\kappa$, say that $S$ has property
$({\rm CI})_\kappa$ if every family of at most $\kappa$ closed left ideals
has nonempty intersection.

\section{Proof intuition}

The gap between the two action conditions is a quantifier gap.  The weak
condition gives one ideal for each pair $(x,y)$; the strong condition needs a
single ideal that works for all $y$ and also lies below any prescribed
principal ideal $St$.

There are two direct ways to close that gap.  If $S$ has a least left ideal,
the same algebraic core lies in every $I_{x,y}$ and every $St$, so no
topology or limiting argument is needed.  More generally, intersect the
closed ideals belonging to a dense set of $y$'s together with $St$.
Property $({\rm CI})_\kappa$ makes this intersection nonempty, and separate
weak continuity plus weak lower semicontinuity extends the inequalities from
the dense set to all of $K$.

Even without either hypothesis, right reversibility lets one satisfy any
\emph{finite} collection of pair inequalities.  Compactness of $K^K$ then
produces a globally nonexpansive pointwise limit map in the closure of every
principal tail.  The missing preservation of weak continuity and
surjectivity pinpoints why this last observation does not settle the full
question.

\section{The least-left-ideal transfer principle}

\begin{theorem}[algebraic upgrade]\label{thm:least}
If $S$ has a least left ideal $L$, every asymptotically nonexpansive action of
$S$ is super asymptotically nonexpansive.  The same assertion holds for
asymptotically $Q$-nonexpansive actions on locally convex spaces.  No topology
or continuity on $S$ is required.
\end{theorem}

\begin{proof}
Fix $x\in K$ and $t\in S$.  For every $y\in K$, asymptotic
nonexpansiveness supplies a left ideal $I_{x,y}$.  Minimality of $L$ gives
$L\subseteq I_{x,y}$ for every $y$.  Moreover, $St$ is a nonempty left ideal,
so $L\subseteq St$.  Thus $L$ itself is supported by $t$ and every $s\in L$
satisfies (2.1) simultaneously for all $y$.  Taking $J_{x,t}=L$ proves the
claim.  The proof for a seminorm $q$ is identical, separately for each $q$.
\end{proof}

\begin{corollary}[all seven source results]\label{cor:seven}
In each of Theorems 3.1, 3.6, 4.1, 4.2, 4.4 and Propositions 3.11, 3.12 of
\cite{MW}, suppose in addition that $S$ has a least left ideal.  Then the
source's super asymptotic or pointwise eventual hypothesis may be replaced by
asymptotic nonexpansiveness (and by asymptotic $Q$-nonexpansiveness in
Theorem 4.4).
\end{corollary}

\begin{proof}
Theorem~\ref{thm:least} supplies the original action hypothesis in every
case.  Apply the corresponding theorem or proposition of \cite{MW} without
altering any of its remaining assumptions.
\end{proof}

\begin{proposition}[a noncompact amenable family]\label{prop:example}
Let $G$ be any infinite discrete amenable group and put
\[
 S=G\times\{0,1\},\qquad (g,i)(h,j)=(gh,\min\{i,j\}).
\]
Then $S$ is a noncompact discrete amenable semitopological semigroup and
$L=G\times\{0\}$ is its least left ideal.  In particular, Theorem 3.1 of
\cite{MW} holds for asymptotically nonexpansive actions of this $S$.
\end{proposition}

\begin{proof}
If a left ideal $I$ contains $(h,j)$, then left multiplication by all
$(g,0)$ shows that $G\times\{0\}\subseteq I$.  The latter set is itself a
left ideal, hence it is least.  The product is discrete and noncompact.  The
projection $S\to G$ and the two-point semilattice factor show amenability;
equivalently, take the product of invariant means.  All left ideals contain
$L$, so $S$ is right reversible.  Corollary~\ref{cor:seven} now applies.
\end{proof}

For the completely explicit instance $G=\mathbb Z$, the semigroup is also
commutative; its principal left ideals are $L$ on layer $0$ and $S$ on layer
$1$.  The verifier enumerates the finite quotients
$\mathbb Z/n\mathbb Z\times\{0,1\}$ and checks the least-ideal statement
exhaustively.

\section{A density/intersection transfer principle}

\begin{theorem}[closed-ideal upgrade]\label{thm:ci}
Let $S$ be a semitopological semigroup, let $K$ be a subset of a locally
convex space $E$, and let $D\subseteq K$ be dense in the original locally
convex topology with $|D|\leq\kappa$, where $\kappa$ is infinite.  Assume:
\begin{enumerate}[label=(\roman*)]
 \item every principal left ideal $St$ is closed;
 \item $S$ has $({\rm CI})_\kappa$;
 \item the action is separately weakly continuous.
\end{enumerate}
Then every asymptotically $Q$-nonexpansive action is super asymptotically
$Q$-nonexpansive.  In a Banach space it is enough that $D$ be norm dense,
and the conclusion applies to ordinary asymptotic nonexpansiveness.
\end{theorem}

\begin{proof}
Fix $x\in K$, $t\in S$, and a continuous seminorm $q$.  For each $y\in D$
choose a left ideal $I_y$ on which
\[
 q(sx-sy)\leq q(x-y).                                      \tag{5.1}
\]
Its closure $\cl I_y$ is a closed left ideal: for $a\in S$, separate
continuity of multiplication gives
$a\cl I_y\subseteq\cl{(aI_y)}\subseteq\cl I_y$.  Moreover (5.1) remains
true for every $s\in\cl I_y$.  Indeed, along a net $s_\alpha\to s$, separate
weak continuity gives $s_\alpha x\to sx$ and $s_\alpha y\to sy$ weakly,
while every continuous seminorm is weakly lower semicontinuous.

The family
\[
 \{St\}\cup\{\cl I_y:y\in D\}
\]
contains at most $\kappa$ closed left ideals.  Hence its intersection $J$ is
nonempty by $({\rm CI})_\kappa$.  It is a left ideal, lies in $St$, and every
$s\in J$ satisfies (5.1) for all $y\in D$.

For arbitrary $y\in K$, choose a net $(y_\alpha)$ in $D$ converging to $y$
in the original topology.  Then $s y_\alpha\to sy$ weakly, and weak lower
semicontinuity gives
\[
 q(sx-sy)\leq\liminf_\alpha q(sx-sy_\alpha)
 \leq\lim_\alpha q(x-y_\alpha)=q(x-y).
\]
Thus $J$ is the ideal required by super asymptotic
$Q$-nonexpansiveness.
\end{proof}

\begin{remark}
The norm-separable case of Proposition 2.7(b) in \cite{MW} is recovered by
using right reversibility for the finite-intersection property, countable
compactness to intersect the countable family, and C-closedness to ensure
that $St$ is closed.  Theorem~\ref{thm:ci} isolates exactly the ideal-theoretic
content needed at an arbitrary density and does not require compactness of
$S$ itself.
\end{remark}

\section{What right reversibility alone still gives}

Write $T_sx=sx$ and give $K^K$ the product of the weak topologies on a weakly
compact $K$.  Thus $K^K$ is compact.

\begin{lemma}[nonexpansive maps in every tail closure]\label{lem:compact}
Let $S$ be right reversible and semitopological, and let its asymptotically
nonexpansive action on a weakly compact $K$ be separately weakly continuous.
For every $t\in S$, the pointwise-weak closure
\[
 \cl{\{T_s:s\in St\}}^{\,K^K}
\]
contains a globally nonexpansive map $p_t:K\to K$.
\end{lemma}

\begin{proof}
Fix a finite $F\subseteq K\times K$.  For each $(x,y)\in F$, close the ideal
$I_{x,y}$.  The sets $\cl{St}$ and $\cl I_{x,y}$ are closed left ideals.
Right reversibility and induction give a point
\[
 r_F\in\cl{St}\cap\bigcap_{(x,y)\in F}\cl I_{x,y}.
\]
Separate weak continuity and weak lower semicontinuity imply
$\|r_Fx-r_Fy\|\leq\|x-y\|$ for every $(x,y)\in F$.  Also
$T_{r_F}$ belongs to the pointwise-weak closure of $\{T_s:s\in St\}$,
because $r_F\in\cl{St}$.

Direct the finite subsets $F$ by inclusion.  Compactness of $K^K$ supplies a
cluster point $p_t$ of $(T_{r_F})$.  For each fixed pair $(x,y)$ the relevant
inequality holds eventually, and weak lower semicontinuity passes it to the
limit.  Hence $\|p_tx-p_ty\|\leq\|x-y\|$ for all $x,y$.  The tail closure is
closed, so it contains $p_t$.
\end{proof}

\section{Verification, limitations, and novelty}

The proofs above use only four independently checkable facts: closures of
left ideals remain left ideals in a semitopological semigroup; continuous
seminorms are weakly lower semicontinuous; $K^K$ is compact by Tychonoff; and
right reversibility gives finite intersections of closed left ideals.  The
included verifier exhaustively checks Proposition~\ref{prop:example} for
finite cyclic quotients and reports the principal left ideals and every left
ideal.

Eight focused upgrade attempts were made.  Counterexample routes based on
commuting interval maps and left-zero elements did not meet the simultaneous
amenability, reversibility, continuity, and asymptotic inequalities.  Direct
adaptations of the source's minimal-set proof failed at the same uniformity
step.  Lemma~\ref{lem:compact} is the strongest unrestricted reduction found,
but a map in a pointwise-weak closure need not be weakly continuous or
surjective.  Those failures prevent the preimage and minimality arguments in
the source from being transferred.  Accordingly, this packet makes no claim
about Question~5.5 without the additional ideal hypotheses.

Bounded searches of the run indexes, the exact question text, and later work
on strong/asymptotically nonexpansive actions found no explicit statement of
Theorems~\ref{thm:least}--\ref{thm:ci} or a full resolution of Question~5.5.
The arguments are elementary enough that they may be folklore; novelty
confidence is therefore moderate.  Mathematical confidence is high for the
two transfer theorems and the compactification lemma.  Human review should
focus on the intended convention for right reversibility and on whether the
least-left-ideal observation has appeared elsewhere under different
terminology.

\begin{thebibliography}{9}
\bibitem{MW}
B. N. Muoi and N.-C. Wong,
\emph{Fixed point theorems of various nonexpansive actions of semitopological
semigroups on weakly/weak* compact convex sets}, arXiv:2006.15393v4 (2022).
\end{thebibliography}

\end{document}
